% ============================================================================
% Fichier  : partie3/P03_C10_normes_professionnelles.tex
% Titre    : Les Normes Professionnelles de l'Audit
% Auteur   : MINKA MI NGUIDJOÏ Thierry Emmanuel
% Version  : 1.0 -- novembre 2026
% Base   : ACSI C10, repris et développé pour INF3226
% ============================================================================
% Sources : IIA, CRIPP 2017 (édition française, dossier Normes) et Normes
% mondiales d'audit interne (2024 ; 5 domaines, 15 principes, 52 normes ;
% applicables depuis le 9 janvier 2025) ; ISACA, ITAF 4e édition (2020),
% numérotation des normes 1000/1200/1400 ; ISO 19011:2018 (voir C05) ;
% ISO/IEC 27007:2020 ; ISO/IEC 27006 : le dossier Normes contient une
% édition scannée (2015) ; la norme est désormais publiée en plusieurs
% parties (27006-1:2024), à vérifier au catalogue ISO ; IAASB, ISA 315
% révisée (2019) et ISAE 3402 (2011) ; INTOSAI, ISSAI 100, 400, 4000 et
% GUID 5100 (2019) ; IDI, manuel d'audit de conformité (2021), dont la
% relation tripartite (ISSAI 400.35) et la règle de référence (ISSAI
% 4000.14). Indépendance déduite des axiomes 4 et 6 d'A0.
% ============================================================================

\chapter{Les Normes Professionnelles de l'Audit}
\label{chap:normes}

\epigraph{%
    « Dire selon quelle norme on a travaillé,\\
    c'est dire au lecteur comment vous contester. »%
}{Adage des réviseurs}

\niveauChapitre{Fondamental}{%
    Maîtrise du Chapitre~\ref{chap:mission} (principes, charte, cycle de
    la mission)%
}

\begin{boxobjectifs}
\begin{itemize}
    \item Situer les cinq familles de normes qui encadrent l'audit des SI,
          et dire laquelle s'applique à une mission donnée.
    \item Présenter les Normes mondiales d'audit interne de l'IIA (2024) et
          leur filiation avec le CRIPP de 2017.
    \item Présenter ITAF (ISACA), puis l'ISO~19011 et les ISO/IEC~27007 et
          27006 pour l'audit des systèmes de management.
    \item Situer les normes de l'audit légal (ISA~315 révisée, ISAE~3402)
          et celles de l'audit public (ISSAI, GUID~5100).
    \item Analyser une situation de conflit d'intérêts, et montrer que
          l'indépendance découle des axiomes~4 et~6 de l'auditologie.
\end{itemize}
\end{boxobjectifs}

\bigskip

% ============================================================================
\section*{Question directrice : selon quelle norme travaille-t-on~?}
% ============================================================================

À Douala, le comité d'audit de la BANQUE SANAGA a trois demandes. Il veut
que Mme~EKAMBI, directrice de l'audit interne, audite la gestion des
changements après la panne de mars~2026. Le commissaire aux comptes, lui,
veut savoir s'il peut s'appuyer sur le rapport que l'intégrateur DIGICORE
AFRIQUE a remis sur ses propres contrôles. Enfin, la COBAC attend une
évaluation du contrôle interne. Trois travaux, trois lecteurs, et autant
de normes.

À Yaoundé, la Chambre des comptes ouvre au même moment un audit de
conformité de l'ANPP. Le magistrat vérificateur, M.~NJOYA, ne suit ni les
normes de l'IIA ni celles du commissaire aux comptes~: il applique les
normes de l'INTOSAI.

\medskip
\textbf{Plan du chapitre.}
\begin{enumerate}
    \item La carte des normes (section~\ref{sec:c13-carte}).
    \item L'audit interne~: les Normes mondiales de l'IIA
          (section~\ref{sec:c13-iia}).
    \item L'audit des SI~: ITAF (section~\ref{sec:c13-itaf}).
    \item Les systèmes de management~: ISO~19011, ISO/IEC~27007 et 27006
          (section~\ref{sec:c13-iso}).
    \item L'audit légal~: ISA~315 et ISAE~3402 (section~\ref{sec:c13-legal}).
    \item L'audit public~: l'INTOSAI (section~\ref{sec:c13-public}).
    \item La méthode du cours et les normes
          (section~\ref{sec:c13-correspondance}).
    \item L'indépendance, lue par l'auditologie
          (section~\ref{sec:c13-deontologie})~; la déontologie est
          développée au Chapitre~\ref{chap:deontologie}.
\end{enumerate}

\separateur

% ============================================================================
\section{La Carte des Normes}
\label{sec:c13-carte}
% ============================================================================

\subsection{Pourquoi des normes~?}

Une norme professionnelle dit ce que l'auditeur doit faire pour que son
travail soit reconnu par ses pairs. Elle a une fonction moins visible~:
elle fixe la règle selon laquelle une conclusion pourra être contestée.
Un rapport qui déclare « nous avons conduit nos travaux conformément
à… » donne au lecteur le moyen de vérifier que l'auditeur a fait ce qu'il
devait. C'est l'axiome d'opposabilité (axiome~6) traduit en pratique
professionnelle~\cite{minka2026auditologie}.

\begin{boxdef}
\termecle{Norme professionnelle d'audit}\index{norme professionnelle
d'audit}~: ensemble d'exigences, émis par une organisation professionnelle
ou normative, qui encadre la conduite d'un audit, la compétence et
l'éthique de l'auditeur, ainsi que la forme de ses conclusions. Elle se
distingue du \emph{référentiel de contrôle} (COSO, COBIT, ISO/IEC~27002),
qui dit \emph{par rapport à quoi} on audite (Chapitre~\ref{chap:referentiels}),
et de la \emph{loi}, qui dit \emph{qui} doit être audité
(Chapitre~\ref{chap:droit}).
\end{boxdef}

\subsection{Cinq familles}

\begin{tableaucours}{L{2.4cm}L{2.3cm}L{3.4cm}L{4.0cm}}{%
    \tableauHeader{Famille} &
    \tableauHeader{Émetteur} &
    \tableauHeader{Textes principaux} &
    \tableauHeader{Dans nos cas}
}
Audit interne & IIA & Normes mondiales (2024)~; CRIPP (2017) & Mme~EKAMBI,
BANQUE SANAGA \\
Audit des SI & ISACA & ITAF (2020)~; Code de déontologie & AUDIX CONSEIL
chez LiZbiZ \\
Systèmes de management & ISO, ISO/IEC & ISO~19011~; ISO/IEC~27007 et
27006 & Audit interne d'un SMSI~; certification \\
Audit légal & IAASB (IFAC) & Normes ISA~; norme ISAE~3402 & Commissaire
aux comptes de la BANQUE SANAGA \\
Audit public & INTOSAI & ISSAI~100, 400, 4000~; GUID~5100 & Chambre des
comptes, audit de l'ANPP \\
\end{tableaucours}
\vspace{-0.8em}
\begin{center}{\small\itshape Tableau~13.1. Les cinq familles de normes de
l'audit des SI}\end{center}

Ces familles ne s'excluent pas. Un auditeur SI certifié par l'ISACA, qui
travaille dans une direction d'audit interne, applique les Normes mondiales
de l'IIA et s'appuie sur ITAF pour les questions techniques. Le choix se
fait d'après le \emph{commanditaire} et le \emph{lecteur} du rapport~: qui
a mandaté la mission, et qui va s'appuyer sur ses conclusions.

% ============================================================================
\section{L'Audit Interne : les Normes Mondiales de l'IIA}
\label{sec:c13-iia}
% ============================================================================

\subsection{Du CRIPP de 2017 aux Normes mondiales de 2024}

Jusqu'en 2024, l'audit interne se référait au \emph{Cadre de référence
international des pratiques professionnelles} (\sigle{CRIPP}), dans son
édition de 2017~\cite{iia2017cripp}. Il comprenait notamment une mission,
des principes fondamentaux, une définition de l'audit interne, un Code de
déontologie en quatre principes et des normes réparties entre normes de
qualification (série~1000) et normes de fonctionnement (série~2000).

En janvier~2024, l'IIA a publié les \textbf{Normes mondiales d'audit
interne} (\textit{Global Internal Audit Standards}), applicables depuis
le 9~janvier~2025~\cite{iia2024gias}. Elles fusionnent dans un seul texte
ce que le CRIPP dispersait, et s'organisent en cinq domaines, quinze
principes et cinquante-deux normes. Chaque norme comporte des exigences,
des considérations de mise en œuvre et des exemples de preuves de
conformité.

\begin{tableaucours}{cL{3.6cm}cL{6.2cm}}{%
    \tableauHeader{Dom.} &
    \tableauHeader{Intitulé} &
    \tableauHeader{Princ.} &
    \tableauHeader{Contenu (notre résumé)}
}
I & Raison d'être de l'audit interne & -- & Ce que l'audit interne apporte
à l'organisation \\
II & Éthique et professionnalisme & 1 à 5 & Intégrité, objectivité,
compétence, conscience professionnelle, confidentialité \\
III & Gouvernance de la fonction d'audit interne & 6 à 8 & Mandat donné
par le conseil, positionnement indépendant, supervision par le conseil \\
IV & Management de la fonction d'audit interne & 9 à 12 & Stratégie,
ressources, communication, qualité \\
V & Réalisation des missions & 13 à 15 & Planification, conduite,
communication des conclusions et suivi des plans d'action \\
\end{tableaucours}
\vspace{-0.8em}
\begin{center}{\small\itshape Tableau~13.2. Les cinq domaines des Normes
mondiales d'audit interne (2024, notre traduction)}\end{center}

\subsection{Ce qui change pour l'auditeur des SI}

\begin{tableaucours}{L{4.6cm}L{7.6cm}}{%
    \tableauHeader{CRIPP 2017} &
    \tableauHeader{Normes mondiales 2024}
}
Code de déontologie, quatre principes & Domaine~II, cinq principes~: la
conscience professionnelle, qui était une norme de qualification, devient
un principe éthique \\
Charte et indépendance (normes 1000 et 1100) & Domaine~III~: le conseil
approuve le mandat, la charte, le plan et les moyens de la fonction \\
Programme d'assurance et d'amélioration de la qualité (normes 1300) &
Principe~12~: évaluation interne continue, évaluation externe au moins
tous les cinq ans \\
Normes de fonctionnement (séries 2000 à 2600) & Domaines~IV et~V,
réorganisés autour de la stratégie de la fonction et du cycle de la
mission \\
\end{tableaucours}
\vspace{-0.8em}
\begin{center}{\small\itshape Tableau~13.3. Filiation entre le CRIPP et les
Normes mondiales (notre lecture)}\end{center}

Les Normes mondiales s'accompagnent d'\emph{exigences thématiques}, qui
précisent ce qu'un auditeur doit couvrir lorsqu'il traite un sujet donné.
La première publiée porte sur la cybersécurité~: l'auditeur des SI d'une
fonction d'audit interne doit s'y reporter.

\begin{boxalert}
\textbf{Violation typique.} La charte d'audit interne de la BANQUE SANAGA,
datée de 2015, renvoie aux « normes de l'IIA » sans préciser l'édition, et
le plan d'audit est approuvé par le directeur général.

\smallskip
\textbf{Règle~:} la charte se réfère aux Normes mondiales de 2024~; le
conseil, ou le comité d'audit qui agit pour lui, approuve la charte, le
plan et les moyens de la fonction (Chapitre~\ref{chap:mission}).
\end{boxalert}

% ============================================================================
\section{L'Audit des SI : ITAF}
\label{sec:c13-itaf}
% ============================================================================

L'ISACA, qui délivre la certification CISA, publie ITAF (\textit{IT Audit
Framework}), cadre de pratiques professionnelles pour l'audit et
l'assurance des SI, dont la quatrième édition date de
2020~\cite{isaca2020itaf}. Ses normes se répartissent en trois séries.

\begin{tableaucours}{L{2.6cm}L{5.0cm}L{4.6cm}}{%
    \tableauHeader{Série} &
    \tableauHeader{Normes (intitulés traduits)} &
    \tableauHeader{Dans ce cours}
}
1000, normes générales & Charte d'audit~; indépendance
organisationnelle~; objectivité~; attente raisonnable~; diligence
professionnelle~; compétence~; assertions~; critères & Chapitre~\ref{chap:mission}~;
fiche d'assertion et critère (Chapitre~\ref{chap:decomposition}) \\
1200, normes de réalisation & Évaluation des risques~; programmation~;
planification de la mission~; réalisation et supervision~; preuves~;
recours aux travaux d'experts~; irrégularités et actes illégaux &
Chapitres~\ref{chap:planifier}, \ref{chap:preuve} et~\ref{chap:fraude} \\
1400, normes de rapport & Rapport~; suivi & Chapitre~\ref{chap:conclure} \\
\end{tableaucours}
\vspace{-0.8em}
\begin{center}{\small\itshape Tableau~13.4. Les trois séries de normes
d'ITAF (4\textsuperscript{e} édition, notre traduction)}\end{center}

Deux normes générales intéressent directement la méthode du cours. La
norme sur les \emph{assertions} demande à l'auditeur d'identifier ce qu'il
vérifie, et la norme sur les \emph{critères} de choisir des critères
pertinents, objectifs et accessibles au lecteur. La fiche d'assertion du
Chapitre~\ref{chap:decomposition} y répond~: chaque contrôle y reçoit un
critère explicite et une question atomique.

\begin{boxinfo}
\textbf{ITAF et l'audit interne.} ITAF ne remplace pas les Normes
mondiales de l'IIA. Un auditeur SI qui travaille dans une direction
d'audit interne applique les Normes mondiales pour l'organisation de la
fonction et le cycle de la mission, et ITAF pour les aspects propres aux
SI~: preuves numériques, recours à des experts techniques, tests sur les
systèmes.
\end{boxinfo}

% ============================================================================
\section{L'Audit des Systèmes de Management}
\label{sec:c13-iso}
% ============================================================================

\subsection{Trois parties, trois audits}

L'ISO~19011:2018, dont le Chapitre~\ref{chap:mission} a présenté les sept
principes, distingue trois situations d'audit selon la relation entre
l'auditeur et l'audité~\cite{iso19011_2018}.

\begin{tableaucours}{L{2.8cm}L{4.6cm}L{4.8cm}}{%
    \tableauHeader{Audit} &
    \tableauHeader{Qui audite qui~?} &
    \tableauHeader{Exemple}
}
Première partie & L'organisation s'audite elle-même (audit interne) &
L'audit interne du SMSI exigé par l'ISO/IEC~27001 (§~9.2) \\
Seconde partie & Un client audite son fournisseur & La BANQUE SANAGA audite
DIGICORE AFRIQUE, en vertu de son contrat \\
Tierce partie & Un organisme indépendant, accrédité, audite pour certifier
& La certification ISO/IEC~27001 d'un centre de données \\
\end{tableaucours}
\vspace{-0.8em}
\begin{center}{\small\itshape Tableau~13.5. Les trois parties de l'audit
des systèmes de management}\end{center}

\subsection{ISO/IEC 27007 et ISO/IEC 27006}

L'\textbf{ISO/IEC~27007:2020} donne des lignes directrices pour auditer un
système de management de la sécurité de l'information (\sigle{SMSI})~:
elle reprend la structure de l'ISO~19011 et ajoute, clause par clause, ce
qui est propre à la sécurité~\cite{iso27007_2020}. Elle s'applique aux
audits de première et de seconde partie, et guide les auditeurs de
certification.

L'\textbf{ISO/IEC~27006} fixe les exigences applicables aux organismes qui
auditent et certifient des SMSI~; elle complète l'ISO/IEC~17021-1, commune
à toutes les certifications de systèmes de management. Elle ne s'adresse
pas à l'auditeur interne, mais elle lui dit ce qu'un certificat garantit, et
ce qu'il ne garantit pas.

\begin{boxinfo}
\textbf{Édition à vérifier.} Le dossier \texttt{Normes} contient une
édition scannée de l'ISO/IEC~27006. La norme est désormais publiée en
plusieurs parties~; avant de la citer dans un rapport, vérifiez l'édition en
vigueur au catalogue de l'ISO, comme pour l'ISO/IEC~27005
(Chapitre~\ref{chap:risque}).
\end{boxinfo}

\begin{boiteRemarque}{Un certificat ISO/IEC 27001 suffit-il~?}
Non. Un certificat atteste qu'un organisme accrédité a jugé le SMSI
conforme, sur un \emph{périmètre} donné, à la date de l'audit. L'auditeur
lit d'abord la déclaration de périmètre et la déclaration d'applicabilité~:
un centre de données certifié peut exclure l'application qui intéresse la
mission. Hors du périmètre, le contrôle reste à coter, et sans autre
preuve, il est \NEK{}.
\end{boiteRemarque}

% ============================================================================
\section{L'Audit Légal : ISA 315 et ISAE 3402}
\label{sec:c13-legal}
% ============================================================================

\subsection{Le commissaire aux comptes et les SI}

Le commissaire aux comptes certifie les états financiers. Ses normes
d'exercice dépendent de l'ordre professionnel national~; les normes
internationales d'audit (\sigle{ISA}) de l'IAASB en sont la référence
mondiale. Son mandat, dans l'espace OHADA, est fixé par
l'AUSCGIE~\cite{ohada2014auscgie} (Chapitre~\ref{chap:droit}).

La norme \textbf{ISA~315}, révisée en 2019, porte sur l'identification et
l'évaluation des risques d'anomalies significatives~\cite{iaasb2019isa315}.
La révision a renforcé ce qu'elle demande sur l'informatique~: l'auditeur
financier doit comprendre l'environnement informatique de l'entité,
identifier les risques qui découlent de l'usage de l'informatique et les
\termecle{contrôles généraux informatiques}\index{contrôles généraux
informatiques} qui y répondent (accès, changements, exploitation). C'est
là que l'auditeur des SI intervient dans la mission légale.

\subsection{Le rapport ISAE 3402}

Lorsqu'une entité confie une partie de ses traitements à un prestataire,
son auditeur ne peut pas toujours aller tester les contrôles chez ce
prestataire. La norme \textbf{ISAE~3402} organise alors un rapport
d'assurance délivré par l'auditeur du prestataire, sur les contrôles de
celui-ci~\cite{iaasb2011isae3402}.

\begin{boxdef}
\termecle{Rapport ISAE~3402}\index{rapport ISAE 3402}~: rapport d'assurance
sur les contrôles d'une société de services, destiné à ses clients et à
leurs auditeurs.
\begin{itemize}
    \item \textbf{Type~1}~: la description du système et la \emph{conception}
          des contrôles, à une date donnée.
    \item \textbf{Type~2}~: la même chose, plus l'\emph{efficacité} des
          contrôles sur une période, avec les tests réalisés et leurs
          résultats.
\end{itemize}
\end{boxdef}

Le rapport énumère aussi les \emph{contrôles complémentaires} que le client
doit mettre en place de son côté~: si le prestataire suppose que le client
revoit chaque trimestre les accès qu'il a demandés, cette revue reste à
tester chez le client.

\begin{boxlecture}{Le rapport de DIGICORE AFRIQUE}
DIGICORE AFRIQUE remet à la BANQUE SANAGA un rapport de type~1, daté du
31~décembre~2025, sur douze contrôles de son service de maintenance du
système central. Le commissaire aux comptes veut conclure sur
l'exercice~2026.

\medskip
\textbf{Constatation.} Le rapport établit la conception des douze contrôles
à une date~; il ne dit rien de leur fonctionnement en 2026.

\medskip
\textbf{Analyse} (confiance élevée). Pour l'efficacité en 2026, les douze
contrôles sont \NEK{}~: le rapport n'est pas faux, il est muet. Le
commissaire aux comptes demande un rapport de type~2 sur la période, ou
réalise ses propres tests. Les trois contrôles complémentaires attendus de
la banque sont à tester chez elle.

\medskip
\hfill\lecture{0{,}00}{1{,}00}{0{,}00}
\end{boxlecture}

Cette lecture rappelle le cas limite du Chapitre~\ref{chap:conclure}~:
une ignorance unanime donne $H = 0$ et $U = 1$. Une entropie nulle n'est
pas une certitude.

\begin{boxinfo}
\textbf{Le commissaire aux comptes.}
\textbf{S'appuyer sur le rapport d'un prestataire}
\begin{enumerate}
    \item Vérifier le type du rapport, la période couverte et le périmètre~:
          les traitements de l'entité y figurent-ils~?
    \item Lire l'opinion de l'auditeur du prestataire et les écarts relevés
          lors de ses tests.
    \item Couvrir la période qui sépare la fin du rapport de la clôture,
          par une lettre de confirmation ou par des tests.
    \item Tester chez l'entité les contrôles complémentaires énumérés dans
          le rapport.
    \item Documenter la conclusion~: ce que le rapport couvre, ce qu'il
          laisse \NEK{}.
\end{enumerate}
\end{boxinfo}

% ============================================================================
\section{L'Audit Public : INTOSAI}
\label{sec:c13-public}
% ============================================================================
\section{La Méthode du Cours et les Normes}
\label{sec:c13-correspondance}
% ============================================================================

La méthode enseignée dans la Partie~\ref{part:methode} n'est pas une norme
de plus~: elle s'inscrit dans les normes existantes. Le Tableau~13.6 montre,
pour chaque étape, l'exigence normative qu'elle met en œuvre. Un rapport
d'audit qui applique la méthode cite la norme du commanditaire, et peut
préciser que la cotation \OUI{}, \NON{}, \NEK{} est la manière dont
l'auditeur documente le caractère suffisant et approprié de ses preuves.

\begin{tableaucours}{L{3.4cm}L{4.6cm}L{4.2cm}}{%
    \tableauHeader{Étape de la méthode} &
    \tableauHeader{Exigence normative} &
    \tableauHeader{Source}
}
Lettre de mission, critères & Objectifs, périmètre et critères définis
avant le terrain & IIA (domaine V)~; ISSAI~4000~; ISO~19011 \\
Décomposition, fiche d'assertion & Identifier l'objet, les critères et les
assertions & ITAF (normes sur les assertions et les critères) \\
Cotation \OUI{}, \NON{}, \NEK{} & Preuves suffisantes et appropriées~;
limitation d'étendue signalée & IIA~; ITAF~; ISSAI~100 \\
Papiers de travail & Documentation permettant à un tiers de refaire le
raisonnement & IIA~; ITAF~; GUID~5100 \\
Contradictoire & Communication des constats et des réponses de la
direction & IIA~; ISSAI~4000 \\
Suivi & Suivi des actions de la direction ou de l'acceptation du risque &
IIA (domaine V) \\
\end{tableaucours}
\vspace{-0.8em}
\begin{center}{\small\itshape Tableau~13.6. La méthode du cours et les
exigences normatives}\end{center}

\begin{boxalert}
\textbf{Violation typique.} Écrire dans un rapport « audit réalisé selon la
méthode \NEK{} », sans citer la norme professionnelle applicable.

\smallskip
\textbf{Règle~:} le rapport cite la norme du commanditaire (IIA, ITAF,
ISO~19011, ISSAI). La méthode du cours est la manière de l'appliquer, pas
un référentiel qui la remplace.
\end{boxalert}

% ============================================================================

\subsection{Le cadre des normes de l'INTOSAI}

Les institutions supérieures de contrôle des finances publiques
(\sigle{ISC}), comme la Chambre des comptes au Cameroun, appliquent le cadre
des prises de position professionnelles de l'INTOSAI. Il comprend trois
niveaux~:
\begin{itemize}
    \item les \textbf{principes} (INTOSAI-P), dont la Déclaration de Lima
          sur les principes du contrôle des finances publiques~;
    \item les \textbf{normes} (\sigle{ISSAI})~: l'ISSAI~100 pose les
          principes fondamentaux et reconnaît trois types d'audit, financier,
          de conformité et de performance~; l'ISSAI~400 énonce les principes
          de l'audit de conformité, et l'ISSAI~4000 en est la norme~;
    \item les \textbf{lignes directrices} (GUID), dont le GUID~5100 sur la
          vérification des systèmes d'information~\cite{intosai2019guid5100}.
\end{itemize}

\subsection{L'audit de conformité}

\begin{boxdef}
\termecle{Audit de conformité}\index{audit de conformité} (ISSAI~400)~:
audit qui évalue si les activités, les transactions et les informations
d'une entité publique sont conformes, dans tous leurs aspects
significatifs, aux autorités qui la régissent (lois, règlements,
décisions, politiques et conditions convenues). Il repose sur une relation
à trois parties~: l'auditeur, la \emph{partie responsable} du sujet et les
\emph{utilisateurs prévus} du rapport~\cite{idi2021conformite}.
\end{boxdef}

Le manuel de l'IDI distingue deux types de missions. Dans une \emph{mission
d'attestation}, l'auditeur se prononce sur une information produite par la
partie responsable. Dans une \emph{mission d'information directe}, il
mesure lui-même le sujet par rapport aux critères. L'assurance donnée peut
être \emph{raisonnable} ou \emph{limitée}~\cite{idi2021conformite}. Les
critères, pertinents, complets, fiables, neutres et compréhensibles, doivent
figurer dans le rapport ou y être référencés~: sans eux, une conclusion
reste ouverte à l'interprétation de chacun.

\begin{boxscenario}{Opération MEFOU. La lettre de la Chambre des comptes}
\textbf{Faits.} M.~NJOYA ouvre un audit de conformité des paiements de
l'ANPP. La partie responsable est le directeur général, M.~ESSOMBA~;
l'utilisateur prévu, le Parlement, destinataire du rapport public. Les
critères sont le régime financier de l'État et les procédures de dépense
de l'agence~\cite{minfi2023regimefinancier}.

\medskip
\textbf{Question de méthode.} Le système de paiement est informatisé. Le
GUID~5100 conduit M.~NJOYA à vérifier d'abord les contrôles généraux
informatiques~: si les comptes des agents partis restent actifs, les
paiements tirés du système ne peuvent être réputés réguliers sans tests
complémentaires. Le Chapitre~\ref{chap:droit} et le ACSI-LAB-06 reprennent ce
cas.
\end{boxscenario}

\begin{boxalert}
\textbf{Violation typique.} Écrire dans un rapport « conformément aux
ISSAI » alors que certaines exigences n'ont pas été suivies.

\smallskip
\textbf{Règle~:} l'ISSAI~4000 interdit de se référer à la norme sans
signaler les exigences non suivies et leurs conséquences~; la même règle
vaut pour toute norme dont on se réclame~\cite{idi2021conformite}.
\end{boxalert}

% ============================================================================
\section{Déontologie et Indépendance}
\label{sec:c13-deontologie}
% ============================================================================

\subsection{Ce que les codes demandent}

Les codes de l'IIA, de l'ISACA et de l'INTOSAI convergent~: intégrité,
objectivité, compétence, conscience professionnelle et confidentialité.
Le Chapitre~\ref{chap:mission} les a présentés à travers les principes de
l'ISO~19011. Le point le plus délicat en pratique est l'indépendance, et
l'atteinte la plus fréquente, le conflit d'intérêts.

\begin{tableaucours}{L{4.4cm}L{3.0cm}L{4.8cm}}{%
    \tableauHeader{Situation} &
    \tableauHeader{Menace} &
    \tableauHeader{Réponse}
}
L'auditrice audite une procédure qu'elle a elle-même conçue & Autorévision
& Déclarer au comité d'audit~; faire évaluer la conception par un tiers \\
L'auditeur est rémunéré au succès de la mission & Intérêt personnel &
Refuser~: les honoraires ne dépendent pas des conclusions \\
L'auditeur audite depuis huit ans le même DSI & Familiarité & Rotation de
l'équipe ou revue par un pair \\
Le directeur général menace de réduire le budget de l'audit & Intimidation
& Saisir le comité d'audit, à qui la fonction rend compte \\
L'auditeur défend publiquement le projet qu'il doit auditer & Défense
d'intérêts & Se retirer de la mission \\
\end{tableaucours}
\vspace{-0.8em}
\begin{center}{\small\itshape Tableau~13.7. Cinq menaces sur
l'indépendance, et la réponse attendue}\end{center}

\subsection{Pourquoi l'indépendance~? La lecture auditologique}

Les normes posent l'indépendance comme une exigence. L'article A0 montre
qu'elle n'est pas une convention~: elle se déduit de deux
axiomes~\cite{minka2026auditologie}.

\begin{boiteTheoreme}{Indépendance (lecture du cours, d'après A0)}
\textbf{Axiome~4 (Interaction stratégique).} L'auditeur et l'audité
poursuivent des intérêts distincts~; l'audité peut retenir l'information,
ou répondre \NEK{} de manière stratégique.

\smallskip
\textbf{Axiome~6 (Opposabilité).} Toute conclusion doit pouvoir être
défendue devant un tiers qui a autorité pour la contester.

\smallskip
\textbf{Conséquence.} Si l'auditeur partage les intérêts de l'audité, le
jeu de l'axiome~4 disparaît~: personne ne conteste l'information retenue.
Le tiers de l'axiome~6 n'a plus alors qu'à montrer cet intérêt commun pour
rendre la conclusion inopposable. L'indépendance est donc une condition
nécessaire de l'opposabilité.
\end{boiteTheoreme}

Cette lecture a une conséquence pratique, qui se mesure. Un \OUI{} donné
par un auditeur en situation de conflit n'est pas opposable~: il ne vaut
pas mieux qu'une réponse inconnue, et se cote \NEK{} jusqu'à ce qu'un tiers
l'ait confirmé.

\begin{boxlecture}{Gestion des changements, SANAGA}
Mme~EKAMBI cote douze contrôles de la gestion des changements~: 9~\OUI{},
1~\NON{}, 2~\NEK{}. Cinq des neuf \OUI{} portent sur la conception de la
procédure qu'elle a rédigée en 2025. Elle le déclare au comité d'audit.

\medskip
\begin{tableaucours}{L{4.4cm}cccccc}{%
    \tableauHeader{Moment} &
    \tableauHeader{\OUI{}} & \tableauHeader{\NON{}} & \tableauHeader{\NEK{}} &
    \tableauHeader{$\mu$} & \tableauHeader{$U$} & \tableauHeader{$H$}
}
Cotation de l'auditrice & 9 & 1 & 2 & 0,75 & 0,17 & 1,04 \\
Cinq \OUI{} non opposables & 4 & 1 & 7 & 0,33 & 0,58 & 1,28 \\
Après revue par un tiers & 7 & 3 & 2 & 0,58 & 0,17 & 1,38 \\
\end{tableaucours}

\medskip
\textbf{Constatation.} La déclaration fait passer le domaine de la
conformité claire à la non-conformité opaque. La revue par un auditeur
externe confirme trois des cinq contrôles et en infirme deux.

\medskip
\textbf{Analyse} (confiance moyenne). La première cotation surestimait la
conformité de 0,17. Le conflit d'intérêts n'est pas une faute morale à
taire~: c'est une source d'incertitude à mesurer, puis à lever.

\medskip
\hfill\lecture{0{,}58}{0{,}17}{1{,}38}
\end{boxlecture}

\begin{boxmetier}{Auditeur SI}
\textbf{Déclarer ses normes et ses limites}
\begin{enumerate}
    \item Indiquer dans la lettre de mission, puis dans le rapport, les
          normes appliquées et leur édition.
    \item Signaler toute exigence non suivie, et sa conséquence sur les
          conclusions.
    \item Déclarer par écrit, avant la mission, toute situation de conflit
          d'intérêts, et la mesure retenue.
    \item Coter \NEK{} un contrôle que l'on ne peut pas évaluer en toute
          indépendance, jusqu'à la revue par un tiers.
    \item Conserver au dossier la déclaration et la revue~: elles font
          partie de la preuve.
\end{enumerate}
\end{boxmetier}

% ============================================================================

\separateur

\begin{boxresume}
\textbf{Ce qu'il faut retenir de ce chapitre~:}
\begin{itemize}
    \item \textbf{Cinq familles de normes}~: audit interne (IIA), audit des
          SI (ISACA), systèmes de management (ISO), audit légal (IAASB),
          audit public (INTOSAI)~; le commanditaire et le lecteur du
          rapport décident de celle qui s'applique.
    \item \textbf{Les Normes mondiales de l'IIA (2024)} remplacent le CRIPP
          de 2017~: cinq domaines, quinze principes, cinquante-deux normes.
    \item \textbf{ITAF} encadre l'audit des SI en trois séries~; ses normes
          sur les assertions et les critères fondent la fiche d'assertion.
    \item \textbf{L'ISO/IEC~27007} guide l'audit d'un SMSI~; l'ISO/IEC~27006
          encadre les organismes de certification~; un certificat ne vaut
          que sur son périmètre.
    \item \textbf{Un rapport ISAE~3402} de type~1 ne dit rien de
          l'efficacité~; seul le type~2 couvre une période.
    \item \textbf{L'audit de conformité public} repose sur trois parties et
          des critères explicites (ISSAI~400 et 4000).
    \item \textbf{L'indépendance} se déduit des axiomes~4 et~6~: un \OUI{}
          donné en situation de conflit se cote \NEK{} jusqu'à la revue par
          un tiers.
    \item \textbf{La méthode du cours} met en œuvre les exigences de ces
          normes~; le rapport cite la norme du commanditaire.
\end{itemize}
\end{boxresume}

\bigskip

\begin{boxexo}[title={\ding{45}\quad Exercice~13.1. Quelle norme~? \niveaubase}]
Pour chaque mission, indiquez la famille de normes applicable et le texte
principal.
\begin{enumerate}[(a)]
    \item La Chambre des comptes vérifie la régularité des paiements de
          l'ANPP.
    \item L'audit interne de la BANQUE SANAGA audite le plan de continuité.
    \item Un organisme accrédité certifie le SMSI d'un centre de données de
          LiZbiZ.
    \item Le commissaire aux comptes de LiZbiZ évalue les contrôles généraux
          informatiques de l'application intégrée.
    \item LiZbiZ audite les contrôles de son hébergeur, en vertu du contrat.
\end{enumerate}
\end{boxexo}

\begin{boxexo}[title={\ding{45}\quad Exercice~13.2. Lire un rapport ISAE 3402 \niveaumoyen}]
DIGICORE AFRIQUE remet en juin~2026 un rapport de type~2 couvrant la
période du 1\textsuperscript{er}~janvier au 31~mars~2026. Sur douze
contrôles, l'auditeur du prestataire relève un écart sur le contrôle des
changements d'urgence. Le rapport attend de la banque trois contrôles
complémentaires.
\begin{enumerate}[(a)]
    \item Quelle partie de l'exercice~2026 le rapport couvre-t-il~?
    \item Comment le commissaire aux comptes couvre-t-il le reste de
          l'exercice~?
    \item Quelle cotation retenir pour le contrôle des changements
          d'urgence, et pour les trois contrôles complémentaires avant tout
          test à la banque~?
    \item Le contrôle des changements d'urgence est au cœur de la panne de
          mars~2026. Quelle conséquence en tirez-vous pour la mission de
          Mme~EKAMBI~?
\end{enumerate}
\end{boxexo}

\begin{boxexo}[title={\ding{45}\quad Exercice~13.3. L'indépendance à l'épreuve \niveauexpert}]
LiZbiZ envisage de confier à AUDIX CONSEIL, qui vient d'auditer son
contrôle d'accès, la refonte de ses profils de droits, puis l'audit de
suivi l'année suivante.
\begin{enumerate}[(a)]
    \item Quelles menaces sur l'indépendance cette double mission
          fait-elle naître~?
    \item À l'aide des axiomes~4 et~6, montrez pourquoi les conclusions de
          l'audit de suivi seraient contestables.
    \item Proposez une organisation qui permette à LiZbiZ de bénéficier de
          la connaissance d'AUDIX CONSEIL sans perdre l'opposabilité du
          suivi.
    \item Comment coteriez-vous, dans l'audit de suivi, les contrôles de
          conception des profils si la mission était malgré tout confiée à
          AUDIX CONSEIL~?
\end{enumerate}
\end{boxexo}

\bigskip

\begin{boxexo}[title={\ding{45}\quad Exercice~13.4. Type 1 ou type 2~? \niveaubase}]
Expliquez en trois phrases à un directeur financier pourquoi un rapport
ISAE~3402 de type~1 ne lui permet pas de conclure que les contrôles de son
prestataire ont fonctionné en 2026.
\end{boxexo}

\begin{boxexo}[title={\ding{45}\quad Exercice~13.5. Un certificat ISO/IEC 27001 \niveaumoyen}]
DIGICORE AFRIQUE présente un certificat ISO/IEC~27001 dont le périmètre est
« hébergement du centre de données de Douala ». La banque lui a aussi
confié l'exploitation de son système central. Que vaut le certificat pour
l'auditrice de la banque~? Quelles questions pose-t-elle~?
\end{boxexo}

\begin{boxexo}[title={\ding{45}\quad Exercice~13.6. Les trois parties \niveaumoyen}]
Pour l'audit de conformité des paiements de l'ANPP, identifiez la partie
responsable, l'auditeur et les utilisateurs prévus. Pourquoi l'ISSAI~400
exige-t-elle que les critères soient connus de la partie responsable~?
\end{boxexo}

\begin{boxexo}[title={\ding{45}\quad Exercice~13.7. Citer la norme \niveaumoyen}]
Rédigez le paragraphe « Normes appliquées » du rapport d'AUDIX CONSEIL chez
LiZbiZ, puis celui du rapport de Mme~EKAMBI à la BANQUE SANAGA.
\end{boxexo}

\begin{boxexo}[title={\ding{45}\quad Exercice~13.8. La correspondance \niveaumoyen}]
À l'aide du Tableau~13.6, montrez qu'un auditeur qui coterait \OUI{} un
contrôle sur la seule déclaration d'un responsable méconnaîtrait au moins
deux exigences normatives. Lesquelles~?
\end{boxexo}

\begin{boxexo}[title={\ding{45}\quad Exercice~13.9. Le conflit à la BANQUE SANAGA \niveauexpert}]
Reprenez la lecture « Gestion des changements, SANAGA » de la
section~\ref{sec:c13-deontologie}. Une partie de l'équipe d'audit interne a
travaillé à la DSI pendant la migration. Proposez une organisation de la
mission qui respecte les Normes mondiales de l'IIA, et dites comment
évoluent les cotes après la revue par un tiers.
\end{boxexo}

\begin{boxexo}[title={\ding{45}\quad Exercice~13.10. Un audit, trois lecteurs \niveauexpert}]
Le comité d'audit de la BANQUE SANAGA veut un seul audit de la gestion des
changements, utilisable à la fois par lui-même, par le commissaire aux
comptes et par la COBAC. Est-ce possible~? Quelles normes, quels livrables,
quelles précautions sur l'indépendance et sur la communication des
résultats~?
\end{boxexo}

\bigskip

\begin{boxinfo}
\textbf{Points de contrôle de l'élève.} Je sais~: nommer les cinq familles
de normes et dire laquelle s'applique~; citer les cinq domaines des Normes
mondiales~; distinguer ISO/IEC~27007 et 27006~; lire un rapport ISAE~3402
de type~1 ou~2~; décrire la relation à trois parties~; relier chaque étape
de la méthode à une exigence normative.
\end{boxinfo}

\begin{boxinfo}
\textbf{Pour aller plus loin.}
\begin{itemize}
    \item IIA, Normes mondiales d'audit interne~\cite{iia2024gias}.
    \item ISACA, ITAF, 4\textsuperscript{e} édition~\cite{isaca2020itaf}.
    \item ISO/IEC~27007:2020~\cite{iso27007_2020}~: l'audit d'un SMSI.
    \item IAASB, ISAE~3402~\cite{iaasb2011isae3402}~: les rapports sur les
          contrôles d'un prestataire.
    \item INTOSAI, GUID~5100~\cite{intosai2019guid5100} et IDI, manuel
          d'audit de conformité~\cite{idi2021conformite}.
\end{itemize}
\end{boxinfo}

\begin{boxinfo}
\textbf{Transition.} Les normes disent comment auditer~; il reste à dire
par rapport à quoi. Le Chapitre~\ref{chap:referentiels} présente les
référentiels de contrôle (COSO, COBIT, ISO/IEC~27002, NIST) et la manière
de les choisir pour une mission.
\end{boxinfo}
